% ================================================================
%  第 12 章 第一部分的注 (原书 p.231-241, 边码 p.141-p.152)
%  源文: _work/src1/task-ch12a.md (12.1–12.22, 注 1–22)
% ================================================================
% ---------------- p.231 ----------------
% 原书 p.231-233 的章内目录(Contents)按体例不排印
\chapter{第一部分的注}

% ---------------- p.232 ----------------
% (章内目录续, 不排印)

% ---------------- p.233 ----------------
\section{$\varnothing$}

\notehead{1}{}
\origpage{141}(1984年3月补写)说我的工作“风格”与“方法”没有变,恐怕失于武断——毕竟我在数学上的表达风格已经深刻改变了。这一年来用于《追寻栈》的大部分时间,是坐在打字机前敲下一些思考;它们就是要几乎原样发表的(只除了后来补加的若干相当简短的注,借互见、勘误之类方便阅读)。没有剪刀,没有浆糊,不必费劲去炮制一份“定稿”(那份定稿尤其不该让任何通向它的过程透出影来)——这毕竟也是“风格”与“方法”的变化!除非把数学工作本身与写作、与陈述结果的工作割裂开来,而那是人为的,因为那不合事情的实情:数学工作与写作是拆不开的。

% ---------------- p.234 ----------------
\section{$\varnothing$}

\notehead{2}{}
\origpage{142}(1984年3月补写)重读最后这两段,我有些不适之感,因为写它们的时候,我牵连的是别人,而不是我自己。显然,写的时候,“我自己也可能与此有关”这个念头并没有碰到我。我肯定什么也没学到:我就这样满足于把多年来在别人身上察知、并以许多方式看着被证实的东西,白纸黑字写下来(无疑还带着某种满意)。在这份反思的后续部分,我不得不回想起:我的一生中,对别人的轻蔑姿态并不缺少。我所把握到的他人之轻蔑与对己之轻蔑之间的关联,若在我这个人身上竟然缺席,那才奇怪;健全的理性(以及类似情形下我对自身视而不见的经验——我终究还是察觉到了)告诉我:肯定不该是这样!然而眼下这只是一条纯粹的推论,它唯一可能的用处,是促使我亲眼看一看究竟发生了什么,去看、去审视(如果它确实存在过、或还存在着)这份仍属假设的对己之轻蔑——它埋得如此之深,至今完全逃过了我的目光。真的,可看的东西并不缺!

\pmark 这一件忽然在我看来是最要紧的几件事之一,恰恰因为它藏得如此之深……\fnote{1}{}{(1984年8月)不过关于这一点,参见“屠杀”一注(n\textsuperscript{o}~87)末两段的反思。}

\section{$\varnothing$}

\notehead{3}{}
我这里想到的尤其是那几条“已故”的猜想——莫德尔(Mordell)的、泰特(Tate)的、沙法列维奇(Chafarévitch)的:去年,这三条猜想在法尔廷斯(Faltings)一篇四十页的手稿里同时得到了证明,而就在那个时候,“圈内人”那套早已稳固的共识还裁定这些猜想“够不着”呢!恰好,我所珍爱的“远阿贝尔代数几何”(géométrie algébrique anabélienne)纲领的那条拱顶石式的基本猜想,就与莫德尔猜想相近。(据说莫德尔猜想甚至是这条基本猜想的推论——这就满可以说明:这个纲领并不是给严肃人预备的故事……)

\section{$\varnothing$}

\notehead{4}{}
再说,即便在今天,也还能碰到地位不明的“证明”。格劳尔特(Grauert)那条以他命名的有限性定理的证明,多年来就是如此:谁(而热心人并不缺!)也读不下去。这种困惑,后来由另一些更透亮的证明解决了;其中有些走得更远,接替了最初那个证明。类似而更极端的情形,是所谓“四色问题”的“解答”:其计算部分是用计算机(外加几百万美元)打发掉的。于是这里所说的“证明”,其根基已不再是理解一种数学处境所生的内心确信,而是对一台没有理解能力的机器的信任——而使用它的数学家,对这台机器的结构与运作一无所知。即便假定这计算由别的计算机、按别的计算程序复核过,我也并不认为四色问题就此了结。它只是换了一副面孔:人们不再去寻一个反例,而只去寻一个证明(当然,得是读得懂的证明!)。

% ---------------- p.235 ----------------
\section{$\varnothing$}

\notehead{5}{}
\origpage{143}\origpage{144}这件事之所以格外值得注意,是因为直到1957年,布尔巴基(Bourbaki)小组里还不止一位成员对我有所保留——而我相信,他们终究把我收编进来时,多少是有些勉强的。一句善意的玩笑话,把我归入“危险的专门家”(指泛函分析方面的)。有时我在嘉当(Cartan)身上感到一种未说出口的、更认真的保留——有那么几年,我一定给了他这样的印象:一个倾向于无端而肤浅地推广的人。我见过他大为惊讶:在我为布尔巴基写的头一篇(也是唯一一篇)稍长的稿子(论流形上的微分体系)里,他发现了一点还算有点内容的思考——当初我提议由我来写时,他并不怎么热心。(这份思考几年后对我又有用:我从凝聚对偶的观点发展留数形式体系时,用到了它。)再说,在布尔巴基大会上,我多半是跟不上的,尤其是在共同宣读稿子的时候:那些宣读与讨论的节奏,我实在跟不下来。也许我真的不适合集体工作。不管怎样,我这种难以融入共同工作的困难,或者我因别的缘故在嘉当等人那里引起的保留,从未在任何时刻给我招来讥讽或冷遇,哪怕一丝居高临下的影子也没有——顶多是韦伊(Weil)那里有过一两次(此人确乎另当别论!)。嘉当对我的善意始终如一,不曾稍减,那里面既有热诚,也有他那特有的、在我看来与他这个人分不开的一点幽默。

\section{《生存与生活》的朋友们}

\notehead{6}{}
在这些朋友里,我大概还该算上皮埃尔·萨米埃尔(Pierre Samuel)——从前我主要是在布尔巴基里认识他的,就像谢瓦莱(Chevalley)一样;他(和谢瓦莱一样)在“生存与生活”(Survivre et Vivre)小组里起了重要作用。我看萨米埃尔并不怎么热衷于那种科学家的优越感的幻觉。我感到,他带来的东西很多,主要靠的是他在共同工作、讨论、与人交往中那份健全的常识和笑呵呵的好脾气,也靠他能在这样一个倾向于激进分析与抉择的群体里,优雅地扮演“讨厌的改良派”的角色。我退出“生存与生活”之后,他还在那里待了一阵,担任同名刊物的主编;当他感到自己在这个群体里的存在已不再有用时,他就体面地离开了(去加入“地球之友”)。

萨米埃尔和我同出那个狭小的圈子;这并不妨碍他属于那些沸腾年月里让我觉得学到了一点东西的朋友之列(尽管我是个多坏的学生……)。他那种为人处世的方式,以及谢瓦莱的方式——两者其实毫不相似——对我那些“精英至上”的倾向来说,是比最犀利的分析更好的解药!

如今在我看来,对那个时期所有让我学到了一点东西的朋友来说,我学到的东西,更多是靠他们各自的为人、靠他们与我不同的感受力(其中“某种东西”终于传达了过来),而不是靠解释、讨论之类。说到这一点,我尤其记得——除了谢瓦莱和萨米埃尔——德尼·盖吉(Denis Guedj)(他在“生存与生活”小组里很有威望)、达尼埃尔·西博尼(Daniel Sibony)(他始终与这个小组保持距离,一边从角落里用半是轻蔑、半是打趣的眼光看着它演变)、戈登·爱德华兹(Gordon Edwards)(他是1970年6月蒙特利尔那场“运动”诞生的共同推动者之一,多年来施展出惊人的精力,维持着《生存与生活》刊物的英文“美洲版”)、让·德洛尔(Jean Delord)(一位年纪与我相仿的物理学家,细致而热忱,他对我连同“生存派”这个小圈子都怀有偏爱)、弗雷德·斯内尔(Fred Snell)(另一位在美国布法罗的物理学家,1972年我在那里住了几个月,做客他的乡间住所)。

% ---------------- p.236 ----------------
\origpage{145}这些朋友当中,五位是数学家,两位是物理学家,而全都是科学家——这似乎表明:那些年里离我最近的环境,仍然是科学家的环境,尤其是数学家的环境。

\section{$\varnothing$}

\notehead{7}{}
上一段是整个导论里第一段在我原稿上被狠狠划掉、并添了许多涂改的。对那件事的描述、对词句的选择,起初都是逆着毛捋、顶着潮流来的——显然有一股力量在推着,要把那件事快快带过,仿佛是了却一桩心事,好“去谈正经事”。这些都是一道抵抗的熟面孔;在这里,它抵抗的是把这段插曲弄清楚,以及弄清它作为内心态度之显影的分量。情形与这篇导论开头(第2段)所描述的完全相似:在数学工作中发现一个矛盾及其意义的那个“关键”时刻——那时充当“抵抗”的,是心智的惰性,是它不肯离开一个错误或不充分的视野(而在那里面,我们这个人丝毫未被牵连)。这种抵抗是主动的,必要时还会发明创造,好把一条鱼在没有水的地方也淹死;而我方才说的惰性,只是一种被动的力。在眼前这桩事上,比在数学工作中尤甚:刚刚以其全部单纯、全部显明浮现出来的那份发现,随即带来的是一种卸下重负的轻松感,一种解放感。这不只是一种感觉——不如说,是对刚刚发生之事的敏锐而心怀感激的知觉,而那件事,就是一次解放。

\section{$\varnothing$}

\notehead{8}{}
后文会看清:这份模棱两可,丝毫没有“在1970年觉醒的次日消散”。这里有一招“自我”典型的战略性撤退:把“觉醒之前”那段时间丢给盈亏自负,而“觉醒”立刻成了一道分界线,划出一个无可指摘的“之后”!

\section{$\varnothing$}

\notehead{9}{}
这不完全确切:在我最亲近的同行里至少有一个例外,后面会看出来。这里有一种记忆特有的“懒惰”:它往往倾向于把那些与一套熟悉且根深蒂固的看法“对不上”的事实“蒙混过去”。

\section{$\varnothing$}

\notehead{10}{}
比如,我写给同事或旧日的学生、那些我当作朋友的人的信——数学问题也好,实际事务或私事也好——多得我已数不清,却从未收到过回音。看来这并不只是专门留给我一个人的“优待”,而是一种风气变化的征兆——我听到的回声也是这个方向。(诚然,那些回声涉及的情形是:寄数学信的人,收信的那位知名数学家并不认识……)

% ---------------- p.237 ----------------
\section{阿尔多·安德烈奥蒂、约内尔·布库尔}

\notehead{11}{}
\origpage{146}\origpage{147}当然。我这边也可能有遗忘——何况我那时格外“硬派”(polar)的性情,恐怕也不会让人乐意跟我谈这类事情,也不会让我记住一场很可能发生过的、这个方向的谈话。可以肯定的是:关于恐惧的问题被提起(甚至不用这个名字来称呼它……),至少该是极其罕见的;今天大概同样罕见,尤其是在“上流社会”里。

在那个圈子里我的众多朋友中,除了谢瓦莱(他至少在六十年代该已察觉到这种恐惧的气氛),我觉得真正清楚地感受到了它的,只有阿尔多·安德烈奥蒂(Aldo Andreotti)。我认识他,认识他妻子芭芭拉(Barbara)和他们那对双胞胎孩子(那时还很小),是在1955年(我想是在芝加哥韦伊家的一次晚会上)。我们一直十分亲近,直到1970年那个“大转折”:我离开了从前属于我们的那个圈子,与他们也就有些失去了联系。阿尔多感受力极敏锐,而这份敏锐丝毫没有因为与数学打交道、与我这样的“硬派”打交道而变钝。他身上有一种对接近他的人的、自然而然的同感之才。这使他与我在这数学圈里、乃至圈外认识的所有别的朋友都不同。在他那里,友情总是压过共同的数学利益(那些利益并不缺),他是极少数与我多少谈过我的一生、他也谈过他一生的数学家之一。他父亲和我父亲一样是犹太人,他因此在墨索里尼的意大利吃过苦头,正如我在希特勒的德国。我见他总是愿意鼓励和支持年轻研究者,而那时的气候,已经越来越难让这些人被建制接纳。他那种自发的关切,总先落在人身上,而不是落在什么数学“潜力”或名声上。他是我有幸遇见的最令人眷爱的人之一。

这番对阿尔多的追忆,又勾起了约内尔·布库尔(Ionel Bucur)的记忆——他也是骤然辞世,未及天年;和阿尔多一样,人们(我想)更怀念的是他这个人、这位乐于重逢的朋友,而不是作为数学讨论伙伴的他。在他身上,人感到一种善良,旁边还有一份少见的谦逊,一种不断让自己退到后面的倾向。这样一个人——如此不把自己当回事、如此不想让谁敬畏——最后竟当上了布加勒斯特大学理学院院长,真是个谜;多半是因为他压根没想到要推辞那些他远不觊觎的职务,而他的同事或政治当局偏要把它们放到他肩上——顺带说,那肩膀是结实的。他是农民的儿子(在一个“阶级标准”要紧的国家里,这一定起了作用),也带着农民那份健全的常识与朴素。他肯定察觉得到围绕着名望之人的那份恐惧,但肯定也觉得这事理所当然,是权力位置的天然属性。不过我并不认为他本人曾让谁感到恐惧——无论是他妻子弗洛丽卡(Florica)、他们的女儿亚历山德拉(Alexandra),还是他的同事或学生;我能听到的回声也确实如此。

\section{$\varnothing$}

\notehead{12}{}
这里的“次日”(lendemain)一词要按字面理解,不是隐喻。

% ---------------- p.238 ----------------
\section{$\varnothing$}

\notehead{13}{}
\origpage{148}很清楚,前面的描述没有别的野心,只是想用具体的词语尽力把记忆的那片“迷雾”所交给我的东西还原出来;这片迷雾从未凝结成任何一个稍为具体的个例,好让我在这里给出某种稍微“写实”或“客观”的描述。若要从这段话里读出:那些不愿坐在前排的同行,或者那些没有明星或泰斗身份的同行,在跟后一类人说话时必定被焦虑缠住——那就是曲解我的意思了。在我于这个圈子里认识的多数朋友身上,显然不是这样,即便是那些偶尔出入讨论会或讨论班的人。毫无保留地成立的是:“泰斗”身份造成一道屏障、一道沟,把它与没有同等身份的人隔开;而这道沟极少消失,哪怕只有一场讨论的工夫。我还要补充:“前排”与“沼泽”之间那条主观的界线(我却觉得它十分真实),绝不能归结为社会学标准(社会地位、职位、头衔等等),甚至也不能归结为“身份”或名声的标准;它还反映着气质上、或者某些更难把握的禀性上的心理特点。我二十岁初到巴黎时,我知道自己是个数学家,我做过数学;尽管有人地生疏之感(我有机会谈到过),我在心底里觉得自己“是他们中的一个”——而只有我知道这一点,甚至一开始我都不确定自己会不会继续做数学。今天,我倒更愿意坐到最后一排去(在难得碰上这个问题的场合)。

\section{$\varnothing$}

\notehead{14}{}
人们可能以为这与“没有首领”的说法相矛盾,其实不然。对布尔巴基的老人们来说,我觉得韦伊被看作小组的灵魂,却从来不是“首领”。他在场、并且他乐意的时候,就成了我所说的“串场人”,但他并不发号施令。他心情不好的时候,会把自己厌恶的某个题目卡住,让别人到下一次大会、趁韦伊不在时再安安静静地重拾这个题目,甚至第二天等他不捣乱了再谈。决议由在场的成员一致通过,因为他们认为:一个人逆着所有其他人的一致意见而站在对的一边,这绝非不可能(甚至并不少见)。这条原则对集体工作来说也许显得荒谬。了不起的是:它居然行得通!

\section{$\varnothing$}

\notehead{15}{}
我并不觉得这种对布尔巴基风格的“过敏”,曾在这些数学家与我、或与布尔巴基其他成员和同情者之间造成交流上的困难——倘若小组的精神真是小圈子精神、是精英中的精英的精神,那倒会如此。在风格与时尚之上,小组的所有成员都有一种对数学实质的敏锐感觉,不论那实质来自何处。只是到了六十年代,我才记得我某位朋友把某些他不感兴趣的数学家称作“讨厌鬼”。这类事我另外几乎一无所知,于是往往把这类评价照单全收,被那份满不在乎的自信震住了——直到有一天我发现,某位“讨厌鬼”是个独到而深刻的心灵,只是没有讨得我那位出色的朋友欢心罢了。我觉得在布尔巴基的某些成员身上,面对别人的工作、在自己不懂或懂得不透的时候所持的谦逊态度(或至少是保留态度),是先被磨掉的;而那种“数学本能”却还留着——那种不必参照名声与声望、就能嗅出丰厚实质或扎实工作的本能。

% ---------------- p.239 ----------------
\origpage{149}从各处传来的回声看,我觉得这两样——谦逊与本能——在从前的那个数学圈子里,今天都已成了稀罕之物。

\section{$\varnothing$}

\notehead{16}{}
说实话,布尔巴基的成员里有好几位肯定有“自己的”小宇宙,大小不一,在布尔巴基小宇宙之外或之上。但就我自己的情形而言,这样一个“自己的宇宙”之所以要等到我不再是布尔巴基的一员、全部精力都投到属于我个人的工作上之后,才在我周围形成——这也许并非偶然。

\section{$\varnothing$}

\notehead{17}{}
主要是在科学圈之外,我才遇到对我所投入的那场行动的热烈回响,以及切实的帮助。除了阿兰·拉斯库(Alain Lascoux)与罗歇·戈德芒(Roger Godement)的友情支持,这里我还尤其要提到让·迪厄多内(Jean Dieudonné)的支持:他特地跑到蒙彼利埃的轻罪法庭去出庭,为一件注定失败的事业,把他那份热忱的证词加在别的证词之上。

\section{$\varnothing$}

\notehead{18}{}
我想,这份缺乏辨别力,在这两个场合并非出于我的疏忽,而是出于不成熟、出于无知。直到十来年后,我才开始注意那些阻滞机制——在我自己身上,也在我的亲人或学生身上——并掂量它们在每个人生活里所起的巨大作用,而且不只是在中学或大学里。当然,我后悔这两个场合没有更成熟一些的辨别力;但我并不后悔把自己的印象清楚地说出来,不论那些印象有没有根据。当我在某个场合看到一件做得不认真的工作,把这事照它的本来面目说出来,我觉得是必要而有益的。若在另一个场合,我得出的结论是仓促的、没有根据的,那么负有责任的也并非只有我一个。被这样摇醒的学生,当时还有得选:要么从中吸取教训(第一次也许正是这样),要么就此泄气,也许还改了行(那也未必是坏事!)。

\section{耶稣与十二使徒}

\notehead{19}{}
从1970年至今,还有一个学生伊夫·拉德加耶里(Yves Ladegaillerie)在我这里准备并通过了论文。第一时期的学生是:P.贝尔特洛(P.~Berthelot)、M.德马祖尔(M.~Demazure)、J.吉罗(J.~Giraud)、M.阿基姆夫人(Mme~M.~Hakim)、黄春生夫人(Mme~Hoang Xuan Sinh)、L.伊利西(L.~Illusie)、P.茹阿诺卢(P.~Jouanolou)、M.雷诺(M.~Raynaud)、M.雷诺夫人(Mme~M.~Raynaud)、N.萨阿韦德拉(N.~Saavedra)、J.-L.维尔迪耶(J.-L.~Verdier)。(其中六位是在1970年之后才完成论文工作的,也就是说,在我数学上的可用时间极为有限的时期。)在这些学生当中,米歇尔·雷诺(Michel Raynaud)占一个特殊位置:他论文工作所处理的那些本质性问题与概念,是他自己找到的,而且他是完全独立地发展它们的;我所谓“导师”的角色,因此仅限于读已完成的论文、组成评审委员会,并成为其中一员。

轮到我出题目的时候,我十分小心,只限于那些我有足够强的关系、因而觉得自己在需要时能撑住学生工作的题目。一个显著的例外,是米歇尔·雷诺夫人(Mme Michèle Raynaud)关于基本群局部与整体莱夫谢茨定理的工作,那些定理是在适当的平展\emph{site}上以 1-栈(1-champs)的措辞表述的。这个问题在我看来是困难的(后来也确实证明如此),而对我所提出的那些猜想,我并没有证明的想法(再说那些猜想也几乎不可能有什么疑问)。这项工作一直做到七十年代初;雷诺夫人(如同此前她丈夫的情形)发展出一套精细而独到的方法,没有得到我或任何别人的帮助。这项出色的工作还引出一个问题:把雷诺夫人的结果推广到 $n$-栈($n$-champs)的情形——在我看来,在概形的语境里,这应当就是“弱莱夫谢茨定理”一类定理的自然归宿。不过,这里相关猜想的表述(它同样几乎不可能有什么疑问)本质上用到了 $n$-栈的概念;而对它的追寻,按说正是本书的主要对象\fnote{2}{}{(其实指的是《数学反思》第3卷,而不是眼前这第1卷《收获与播种》——见导论,第(v)页。)},正如书名《追寻栈》所示。我们大概会在相应的地方再回到这一点。

% ---------------- p.240 ----------------
\origpage{151}\origpage{150}另一个相当特殊的情形是黄春生夫人。我最初是在1967年12月在河内遇见她的,那时我在疏散后的河内大学讲了一个月的课程—讨论班。第二年,我向她提出了她的论文题目。她是在战争时期格外艰难的条件下工作的,与我的接触仅限于断续的通信。1974/75年她得以来到法国(趁温哥华国际数学家大会之便),于是在巴黎通过了论文(评审委员会由嘉当主持,另有施瓦茨(Schwartz)、德尼(Deny)、齐斯曼(Zisman)和我)。

最后我还该提到皮埃尔·德利涅(Pierre Deligne)和卡洛斯·孔图-卡雷尔(Carlos Contou-Carrère):这两人都多少算是学生,前者在1965—68年前后,后者在1974—76年前后。两人显然都有(并且一直有)非同寻常的才能,而他们运用这才能的方式极为不同,际遇也极为不同。到比雷(Bures)来之前,德利涅在比利时多少算是蒂茨(Tits)的学生——不过我怀疑,按通常意义说,他在数学上做过谁的学生。孔图-卡雷尔做过桑塔洛(Santalo)的学生(在阿根廷),有一阵还做过托姆(Thom)的学生(多少算吧!)。接触建立的时候,两人都已具备一个数学家的身量,只差孔图-卡雷尔缺少方法与行当上的训练。

我在数学上对德利涅的作用,仅限于零零碎碎地把我所懂的那点代数几何讲给他听;他学起来就像听故事——仿佛他向来就知道。顺带我也提出一些问题,而他多半当场或在随后几天里就找到了答案。这就是我所知道的德利涅最早的工作。1970年以后的工作(他的,以及我那些“正式学生”的),我只从极为零散遥远的回声里略知一二\fnote{3}{}{(我尤其有机会翻阅过贝尔特洛和德利涅寄给我的几份抽印本,承他们的好意。)}。

至于我对孔图-卡雷尔的作用,按他自己在论文开头所说,仅限于把他引进概形的语言。无论如何,他这几年作为国家博士论文所准备的工作,我只是远远地跟着;那是一个再时新不过的题目,超出了我的能力。

% ---------------- p.241 ----------------
\origpage{152}孔图-卡雷尔在大世界里经历了几番波折之后,最近终于在最后关头——而且(现在在我看来)是违着他的心意——不得不求我帮忙,充当论文导师并组成评审委员会。(这有让他显得像“1970年以后”的格罗滕迪克学生的风险,而在一种处境里,这可能会带来严重的不便……)我尽力履行了这项任务;这大概是我最后一次行使这个职能(在国家博士论文这一级上)。在这个有点特殊的场合,我格外高兴有让·吉罗(Jean Giraud)的友情协助:他也拨出一两个月的时间,仔细读了那部厚厚的稿子,并写了一份详尽而热忱的报告。

\section{$\varnothing$}

\notehead{20}{}
这让我想起莫妮克·阿基姆(Monique Hakim)所选的题目,说实话那也不见得更有吸引力;我真不知道她是怎么保持住精神的!如果说她有时也吃力,那至少没有到让她郁郁寡欢、闷闷不乐的地步;我们之间的工作是在亲切而放松的气氛里进行的。

\section{$\varnothing$}

\notehead{21}{}
也许更确切的说法是:对我这样的气质来说,要完全担起一个教师角色,我缺的仍是那份必要的成熟。我后天形成的气质,长期带着“男性”(或“阳”)特征过度占优的印记;而成熟的一个方面,恰恰是一种以“女性”(或“阴”)为主的“阴阳”平衡。

(后来补写)我看出,比成熟更缺的,是一种慷慨——在我至今的教学生涯里,我缺的正是某种慷慨;这种慷慨的表现方式,比在时间与精力上随时可用更为细微,也更为本质。在我第一个教学时期,这份欠缺并没有明显地表现出来(比如说,通过一连串失败的情形积累起来),多半尤其是因为它被选择来跟我工作的学生身上强烈的动机补偿了。相反,在第二个时期,从1970年至今,我觉得这份欠缺至少是原因之一,而且无论如何是最直接牵涉到我自己的那一条原因——我在研究层面(也就是从DEA这一级起)的教学中所看到的整体失败,原因就在这里。关于这一点,见《纲领概要》(Esquisse d'un programme)第8段,以及第9段“一项教学活动的总结”;在那里面透出的是这七八年来这项活动留给我的挫败感\fnote{4}{}{(另可参见后来补写的注(23iv)。)}。

\section{$\varnothing$}

\notehead{22}{}
也许不会再有多久了,因为我已经决定申请进入国家科学研究中心(CNRS),从而结束大学里的教学活动——这几年来,这活动变得越来越成问题。

% ---------------- p.242 ----------------
\origpage{153}\origpage{154}
\section{$\varnothing$}

\notehead{22$'$}{}
即便在1970年以后,当我对数学的兴趣在我生活里变得时有时无、退到了边缘,我也不认为有过哪一回,学生来叫我同他一起工作、而我推辞掉的。我甚至可以说:除了两三个例外,我1970年以后的学生对自己所做工作的兴趣,远远低于我本人对他们题目的兴趣——即便在我除了去系里露面的日子几乎不碰数学的那些时期,也是如此。因此,我1970年以前对学生的那份随时可供支配,以及作为其首要标志的、工作上的极端严格,对我后来多数学生来说恐怕毫无意义:他们做数学并没有什么信念,仿佛那是一种他们不得不对自己持续施加的努力……

\section{孩子与老师}

\notehead{23}{}
这里“传递”一词,与事情的实情并不真正相符——实情把我拉回到一种更谦逊的态度上来。这种严谨并不是一件可以传递的东西,最多只能被唤醒、被鼓励,而它从小就被忽视、被打消,被家庭环境,也被学校和大学。就我所能回想起来的而言,这种严谨一直存在于我的种种追寻之中——至少是那些智力性质的追寻——而我不认为它是由我父母传给我的,更不认为它传自哪位师长,无论是在学校里,还是在数学界的前辈当中。我觉得它似乎属于“天真”的诸属性之一,因而也属于每个人生来便被分得的东西。这种天真很早就“吃够了苦头”,于是不得不或深或浅地潜下去,以致在往后的一生里往往再也见不到它多少痕迹了。在我身上,由于一些我还不曾想去深究的原因,某种天真在智力好奇心这个相对无害的层面上存活了下来;而在其余一切地方,它都深深地潜了下去,神不知、鬼不觉!就像在每个人身上那样。也许,完全意义上的“教学”的秘密——或者不如说它的奥秘——就在于重新与这种表面看来已经消失的天真接上联系。但是,如果这联系尚未首先存在于、或尚未被重新找回到教师本人身上,那就根本谈不上在学生身上把它重新找到。而这时教师“传给”学生的,绝不是这种严谨或这种天真(两者在师生身上都是天生的),而是一种尊重,一种对这件通常遭人弃置的东西的、无声的重新肯定。

\section{$\varnothing$}

\notehead{23$'$}{}
不过,七八年来,在我的数学生涯里还有另一个慢性的“挫败之源”,只是这些年来它表现得隐晦得多。靠着反复出现,靠着同一种“令人挫败”的情形在我的教学活动中顽固地累积,它终究变得显眼起来,末了还爆发为一声“受够了!”,使我实际上终止了一切所谓的“指导研究工作”。在我的这番反思途中,我有一两次顺带触及这个问题,最后总算在结尾处多少把它审视了一番。我在那里至少描述了这份挫败感,并考察了它在我“重返数学”一事中所起的作用(参见第50段“一段过去的重量”)。
